\documentclass{ximera}

\input{../preamble.tex}


\title[Problem 5]{Problem 5}

\begin{document}
\begin{abstract} \end{abstract}
\maketitle


% Extracted from lHopitalsRule.tex, problem #5
\begin{problem}
Let $f$ and $g$ be two functions such that $\displaystyle \lim_{x\to\infty}f(x)=\displaystyle \lim_{x\to\infty}g(x)=\infty$.\\[1em]

We say that $f$   \textbf{grows faster} than $g$ as $x$ goes to $\infty$  if  $\displaystyle \lim_{x\to\infty}\dfrac{f(x)}{g(x)}=\infty$, or, equivalently, if\\

 $\displaystyle \lim_{x\to\infty}\dfrac{g(x)}{f(x)}=0$.\\

If the limit exists, namely, if  $\displaystyle \lim_{x\to\infty}\dfrac{f(x)}{g(x)}=M$, for some positive number $M$,\\

 then we say that the functions $f$ and $g$ have \textbf{comparable growth rates}.\\

In other words, we compare the growth rates of  functions $f$ and $g$ by computing the limit of the form $\dfrac{\infty}{\infty}$: $\displaystyle \lim_{x\to\infty}\dfrac{f(x)}{g(x)}$ .\\

For each of the following pairs of functions, determine which of the pair grows faster or state that they have comparable growth rates.  Justify your answer by computing the limit $\displaystyle \lim_{x\to\infty}\dfrac{f(x)}{g(x)}$.



\begin{enumerate}
	\item$b^x; x^x; b>1$
		\begin{explanation}
		$\displaystyle \lim_{x \to \infty}\dfrac{b^x}{x^x}=\displaystyle \lim_{x \to \infty}\left(\dfrac{b}{x}\right)^x=0$.  Since this limit is of the form $0^{\infty}$, which is a determinate form.
		Therefore, $x^x$ grows faster.
		
		\end{explanation}
	
	\item $x^x; \left(\dfrac{x}{e}\right)^x$
		\begin{explanation}
		$\displaystyle \lim_{x \to \infty}\dfrac{x^x}{(x/e)^x}=\displaystyle \lim_{x \to \infty}\dfrac{x^x}{\dfrac{x^x}{e^x}}=\displaystyle \lim_{x \to \infty}{x^x}\cdot{\dfrac{e^x}{x^x}}=\displaystyle \lim_{x \to \infty}{e^x}=\infty$  Therefore, $x^x$ grows faster.
		
		\end{explanation}
	

\item $x^3;  x^3 \cdot \ln(x)$
		\begin{explanation}
		$\displaystyle \lim_{x \to \infty}\dfrac{x^3}{x^3 \cdot \ln(x)}=\displaystyle \lim_{x \to \infty}\dfrac{1}{\ln(x)}=0$.  
		Therefore, $x^3 \cdot \ln(x)$ grows faster.
		
		\end{explanation}
	\item $a^x; b^x; 0<a<b$
	\begin{explanation}
		$\displaystyle \lim_{x \to \infty}\dfrac{a^x}{b^x}=\displaystyle \lim_{x \to \infty}\left(\dfrac{a}{b}\right)^x=0$
		Therefore, $b^x$ grows faster
	
	
	\end{explanation}
	\item $\log_a(x); \log_b(x); 1<a<b$
	\begin{explanation}
	$\displaystyle \lim_{x \to \infty}\dfrac{\log_a(x)}{\log_b(x)}$\\
	By L'Hopital's Rule: $\displaystyle \lim_{x \to \infty}\dfrac{\dfrac{1}{x\ln(a)}}{\dfrac{1}{x\ln(b)}}=\displaystyle \lim_{x \to \infty}\left(\dfrac{1}{x\ln(a)} \cdot \dfrac{x\ln(b)}{1}\right)=\dfrac{\ln b}{\ln a}$\\
	Therefore, $\log_a(x)$ and $\log_b(x)$ grow at comparable rates.
	
		\end{explanation}
	\item $\ln^{3} (x); x^{1/2}$
		\begin{explanation}
	$\displaystyle \lim_{x \to \infty}\dfrac{\ln^{3}(x)}{x^{1/2}}$\\
	By L'Hopital's Rule:$\displaystyle \lim_{x \to \infty}\dfrac{\ln^{3}(x)}{x^{1/2}}=\displaystyle \lim_{x \to \infty}\dfrac{3\ln^2 (x) \cdot (1/x)}{(1/2)x^{-1/2}}=\displaystyle \lim_{x \to \infty}\dfrac{6 \cdot \ln^2 (x)}{x^{1/2}}$\\
	By L'Hopital's Rule:$\displaystyle \lim_{x \to \infty}\dfrac{12\ln (x) \cdot (1/x)}{(1/2)x^{-1/2}}=\displaystyle \lim_{x \to \infty}\dfrac{24 \cdot \ln (x)}{x^{1/2}}$\\
	By L'Hopital's Rule: $\displaystyle \lim_{x \to \infty}\dfrac{48}{x^{1/2}}=0$
	Therefore, $x^{1/2}$ grows faster
		\end{explanation}

%part g	
	\item $x; \ln(x)\sqrt{x}$
		\begin{explanation}
	
	$\displaystyle \lim_{x \to \infty}\dfrac{x}{\ln (x)\sqrt{x}}=\displaystyle \lim_{x \to \infty}\dfrac{x^{1/2}}{\ln(x)}$\\
	By L'Hopital's Rule:$\displaystyle \lim_{x \to \infty}\dfrac{(1/2)x^{-1/2}}{x^{-1}}=\displaystyle \lim_{x \to \infty}{(1/2)x^{1/2}}=\infty$\\
		Therefore, $x$ grows faster
		\end{explanation}
		
		%part h
		\item Challenge: $x^{40}; 1.004^x$ (Hint: Use the substitution $x=\ln(t)$.)
			\begin{explanation}
			\begin{align*}
			\displaystyle \lim_{x \to \infty}\dfrac{x^{40}}{1.004^x} &=\displaystyle \lim_{t \to \infty}\dfrac{[\ln(t)]^{40}}{1.004^{\ln(t)}}\\
			&=\displaystyle \lim_{t \to \infty}\dfrac{\ln(t)^{40}}{{e^{\ln(1.004)}}^{\ln(t)}}\\
			&=\displaystyle \lim_{t \to \infty}\dfrac{\ln(t)^{40}}{t^{\ln(1.004)}}\\
			&=\displaystyle \lim_{t \to \infty}\left(\dfrac{\ln(t)}{t^{1.004/40}}\right)^{40}\\
			&=\left[\displaystyle \lim_{t \to \infty}\left(\dfrac{\ln(t)}{t^{1.004/40}}\right)\right]^{40}\\
			&\stackrel{L.R.}{=} \left[\displaystyle \lim_{t \to \infty}\dfrac{\dfrac{1}{t}}{(1.004/40)t^{(1.004/40)-1}}\right]^{40}\\
			&=\left[\displaystyle \lim_{t \to \infty}\dfrac{1}{(1.004/40)t^{((1.004/40)-1+1)}}\right]^{40}\\
			&=0^{40}\\
			&=0	
				\end{align*}
		Therefore, $1.004^x$ grows faster.
		\end{explanation}
	%part i
	\item Put the following functions in order of growth rate.\\
	$$x^3 \cdot \ln(x), \ln^{3} (x), x^x, 1.004^x, x^{40}, x^3$$
	
	\begin{explanation}
	
	$$ \ln^{3} (x)<<x^{3}<<x^3 \cdot \ln(x)<< x^{40}<< 1.004^x <<x^x$$
	
	\end{explanation}
	
\end{enumerate}
\end{problem}



\end{document}
